\subsection{Integration of functions with values in a Banach space, with respect to an integral of point measures}

THEOREM 2. --- Let \((\pi, g)\) be a \(\mu\) -adapted pair, and let

\[
\nu = \int g (t) \varepsilon_ {\pi (t)} d \mu (t).
\]

Let \(\mathbf{f}\) be a function defined on \(\mathbf{X}\) , with values in a Banach space \(\mathbf{F}\) or in \(\overline{\mathbf{R}}\) . For \(\mathbf{f}\) to be essentially \(\nu\) -integrable, it is necessary and sufficient that \(t \mapsto \mathbf{f}(\pi(t))g(t)\) be essentially \(\mu\) -integrable, in which case

\[
\int \mathbf {f} (x) d \nu (x) = \int \mathbf {f} (\pi (t)) g (t) d \mu (t).\tag{9}
\]

Suppose, moreover, that \(\pi\) is continuous and proper, and that g is continuous and such that \(g(t) > 0\) for all \(t \in T\) . Then, for f to be \(\nu\) -integrable, it is necessary and sufficient that \(t \mapsto \mathbf{f}(\pi(t))g(t)\) be \(\mu\) -integrable.

\begin{enumerate}
\def\labelenumi{\Alph{enumi})}
\tightlist
\item
  We begin by treating the case that the measure \(\mu\) has compact support K, on which g is bounded. The measures \(\mu\) and \(\nu\) are then bounded, and one can replace `essentially integrable' in the statement by `integrable'. Suppose that f is \(\nu\) -integrable: the function \(\mathbf{f}(\pi(t))g(t)\) is then \(\mu\) -integrable, and the relation (9) is verified, by Th. 1 of §3, No.~3. Conversely, suppose that \(\mathbf{f}(\pi(t))g(t)\) is \(\mu\) -integrable: f is then \(\nu\) -measurable (No.~3, Prop. 3 and Remark), and
\end{enumerate}

\[
\int^ {\bullet} | \mathbf {f} (x) | d \nu (x) = \int^ {\bullet} | \mathbf {f} (\pi (t)) | g (t) d \mu (t) <   + \infty
\]

(No.~2, Th. 1); \(\mathbf{f}\) is therefore essentially \(\nu\) -integrable (§1, No.~3, Prop. 9), hence \(\nu\) -integrable. Th. 1 of §3, No.~3 then implies (9).

\begin{enumerate}
\def\labelenumi{\Alph{enumi})}
\setcounter{enumi}{1}
\tightlist
\item
  Let us pass to the general case. Let \(\mathfrak{K}\) be the set of compact subsets \(K\) of \(T\) such that \(g|K\) is continuous: \(\mathfrak{K}\) is \(\mu\) -dense (Ch. IV, §5, No.~10, Prop. 15), therefore the measure \(\mu\) is the sum of a family \((\mu_{\alpha})_{\alpha \in A}\) of measures whose supports are elements of \(\mathfrak{K}\) (§2, No.~3, Prop. 4). The pair \((g,\pi)\) is obviously \(\mu_{\alpha}\) -adapted for every \(\alpha \in A\) , and the measure \(\nu\) is the sum of the family of measures \(\nu_{\alpha} = \int \varepsilon_{\pi(t)} g(t) d\mu_{\alpha}(t)\) (§3, No.~1, Cor. of Prop. 1). Since the argument of A) may be applied to the measures \(\mu_{\alpha}, \nu_{\alpha}\) , the first part of the statement then follows from Prop. 3 of §2, No.~2.
\end{enumerate}

For the function \(\mathbf{f}\) (resp. \(t\mapsto \mathbf{f}(\pi (t))g(t)\) ) to be integrable for \(\nu\) (resp. for \(\mu\) ), it is necessary and sufficient that it be essentially integrable and that

\[
\int^ {*} | \mathbf {f} (x) | d \nu (x) <   + \infty \quad (\mathrm{resp.} \int^ {*} | \mathbf {f} (\pi (t)) | g (t) d \mu (t) <   + \infty).
\]

The second part of the statement therefore follows from the first part and Proposition 2.

Remark. --- Let \((\pi, g)\) be a \(\mu\) -adapted pair, \(\pi'\) a mapping of T into X, and \(g'\) a finite numerical function \(\geqslant 0\) defined on T, such that \(\pi'\) (resp. \(g'\) ) is equal to \(\pi\) (resp. g) locally almost everywhere for \(\mu\) . Then the pair \((\pi', g')\) is \(\mu\) -adapted, the measures \(\lambda_t = g(t)\varepsilon_{\pi(t)}\) and \(\lambda_t' = g'(t)\varepsilon_{\pi'(t)}\) are equal locally almost everywhere, and \(\int g(t)\varepsilon_{\pi(t)} d\mu(t) = \int g'(t)\varepsilon_{\pi'(t)} d\mu(t)\) . If now \(\pi'\) and \(g'\) are only defined locally almost everywhere (for \(\mu\) ) and if there exists a \(\mu\) -adapted pair \((\pi, g)\) such that \(\pi'\) (resp. \(g'\) ) is equal to \(\pi\) (resp. g) locally almost everywhere, one again says that the pair \((\pi', g')\) is \(\mu\) -adapted and one sets

\[
\int g ^ {\prime} (t) \varepsilon_ {\pi^ {\prime} (t)} d \mu (t) = \int g (t) \varepsilon_ {\pi (t)} d \mu (t)
\]

(cf.~§3, No.~3, Remark). The statements of Ths. 1 and 2 and of Prop. 3 remain valid when \(\pi\) and \(g\) are only assumed to be defined locally almost everywhere.

